\documentclass[10pt,a4paper]{amsart}
\usepackage[margin=19mm]{geometry}
\usepackage{amsmath,amssymb,mathtools}
\usepackage[scr=rsfs,cal=euler]{mathalfa}
\usepackage{xfrac}
\usepackage[unicode,hidelinks,bookmarks=false]{hyperref}
\newcommand{\ie}{i.e.\ }
\newcommand{\cf}{cf.\ }
\newcommand{\resp}{respectively\ }
\newcommand{\Subsection}{\subsection}
\providecommand{\nobreakdash}{\nobreak\hbox{-}}
\setlength{\emergencystretch}{2em}
\multlinegap0pt
\usepackage{enumitem}
\usepackage{bm}
\usepackage{algpseudocode}
\usepackage{algorithm}
\newcommand{\absrajcode}[1]{\mathrm{raj}(#1)} 
\newcommand{\absoverrajcode}[1]{\overline{\mathrm{raj}}(#1)} 
\newcommand{\absinvcode}[1]{\mathrm{inv}(#1)}
\newcommand{\absoverinvcode}[1]{\overline{\mathrm{inv}}(#1)}

\DeclareMathOperator{\nested}{\mathrm{nested}}
\DeclareMathOperator{\greedy}{\mathrm{greedy}}
\DeclareMathOperator{\leap}{\mathrm{leap}}
\DeclareMathOperator{\stair}{\mathrm{stair}}
\DeclareMathOperator{\leads}{\mathrm{leads}}
\DeclareMathOperator{\supp}{\mathrm{supp}}
\DeclareMathOperator{\LIS}{\mathrm{LIS}}
\DeclareMathOperator{\wt}{\mathrm{wt}}
\DeclareMathOperator{\dwt}{\overline{\mathrm{wt}}}
\DeclareMathOperator{\karola}{\mathrm{karola}}
\DeclareMathOperator{\cc}{\mathrm{invcode}} 
\newcommand{\occ}{\overline{\cc}} 
\DeclareMathOperator{\rr}{\mathrm{rajcode}} 
\newcommand{\orr}{\overline{\rr}} 
\newcommand{\LL}{\overline{L}}
\newcommand{\CG}{C^G} 
\newcommand{\CN}{C^N} 
\newcommand{\CL}{C^L} 
\newcommand{\CS}{C^S}

\newtheorem{theorem}{Theorem}[section]
\newtheorem{lemma}[theorem]{Lemma}
\newtheorem{corollary}[theorem]{Corollary}
\newtheorem{proposition}[theorem]{Proposition}
\newtheorem{definition}[theorem]{Definition}
\newtheorem{example}[theorem]{Example}
\newtheorem{remark}[theorem]{Remark}
\newtheorem{conjecture}[theorem]{Conjecture}
\newtheorem{question}[theorem]{Question}
\title[Leading monomials above minimal degree]{On the Degree of Grothendieck Polynomials: original Theorem 11.3}
\author{Matt Dreyer, Karola M\'esz\'aros and Avery St. Dizier}
\date{Source: Algebraic Combinatorics 7(3) (2024), 627--658; DOI 10.5802/alco.358}
\begin{document}\maketitle
\noindent\textit{Editorial scope.} The counted unit is original Theorem 11.3 with its complete authored proof. Because that proof consumes the paper's climbing-chain machinery, the necessary local core is retained as uncounted context: the Schubert, Grothendieck, Rothe, Lehmer, increasing-subsequence and Rajchgot definitions, the climbing-chain, marking and weight definitions with the Lenart-Robinson-Sottile formula, the greedy chain, Algorithm 1 and every lemma and theorem the proof invokes, and Lemmas 11.1 and 11.2 with their proofs. Later staircase and leaping chain sections and the two closing conjectures are not selected. No graphics file is redistributed; the two figure pointers inside the retained proofs resolve by hyperlink to the published Figures 4 and 7 and no case is decided by a picture. The duplicated dangling clause in Lemma 9.3 and the delimiter macros of equation (1) are kept exactly as printed, without generated corrections.
\section{Background: Schubert, Grothendieck and Rajchgot code}
\setcounter{section}{1}\setcounter{theorem}{0}
\begin{theorem}[{\cite[Theorem 1.1]{PSW}}]
	\label{thm:pswleadingterm}
	For any $w\in S_n$,
	\[\deg \mathfrak{G}_w=\absrajcode{w}.\]
	Moreover in any term order satisfying $x_1<x_2<\cdots<x_n$, the leading monomial of the highest-degree homogeneous component $\mathfrak{G}_w^{\mathrm{top}}$ of $\mathfrak{G}_w$ is $\bm{x}^{\rr(w)}$.
\end{theorem}

\setcounter{section}{1}\setcounter{theorem}{1}
\begin{conjecture}[{\cite[Conjecture 4.1]{elena1}}]
	\label{conj:elena}
	Fix $w\in S_n$ and any term order with $x_1<x_2<\cdots<x_n$. For $\ell(w)< k\leq \absrajcode{w}$, let $m_k(\bm{x})$ be the leading monomial of the degree~$k$ homogeneous component of $\mathfrak{G}_w$. Then
	\[m_k(\bm{x}) = x_pm_{k-1}(\bm{x}) \]
	where~$p$ is the largest index such that $x_p m_{k-1}(\bm{x})$ divides $\bm{x}^{\rr(w)}$.
\end{conjecture}

For $n\in \mathbb{N}$, we write $[n]$ to mean the set $\{1,2,\ldots,n \}$. For two indices $1\leq i<j\leq n$ we write $t_{ij}$ for the transposition in the symmetric group $S_n$ swapping~$i$ and~$j$. We write $s_j$ for the adjacent transposition $t_{j,j+1}$.

Throughout, we write permutations in one-line notation (as a string) $w=w(1)w(2)\cdots w(n)$. We will take permutations as acting on the right -- switching positions, not values. For example $ws_1$ equals~$w$ with the numbers $w(1)$ and $w(2)$ swapped. 

For $v\in \mathbb{Z}^n$, we write $|v|$ for $v_1+v_2+\cdots+v_n$. We denote the standard basis vectors of $\mathbb{Z}^n$ by $e_1,e_2,\ldots,e_n$. 
We use the notation $\overline{v}$ to denote the \emph{vector complement} 
\[\overline{v}_k = n-k - v_k \mbox{ for each } k\in[n].\]
A \emph{term order} on $\mathbb{Z}[x_1,\ldots,x_n]$ is a total order~$<$ of all (monic) monomials such that
\[x^{\alpha}\leq x^{\beta}\mbox{ implies } x^{\alpha+\gamma}\leq x^{\beta+\gamma} \mbox{ for all }\gamma\in\mathbb{Z}^n_{\geq 0}. \]
The \emph{leading term} of a polynomial $f\in \mathbb{Z}[x_1,\ldots,x_n]$ is the term of the largest monomial under~$<$ appearing in~$f$. The \emph{leading monomial} is the leading term divided by its coefficient. For example if $f(x_1,x_2)=x_1^2+2x_1x_2+3x_2^2$ and $x_1<x_2$, the leading term of~$f$ is $3x_2^2$ and the leading monomial is $x_2^2$.

\setcounter{section}{2}\setcounter{theorem}{0}
\begin{definition}
	Fix any $n\geq 0$. The \emph{divided difference operators} $\partial_j$ for $j\in[n-1]$ are operators on the polynomial ring $\mathbb{Z}[x_1,\ldots,x_n]$ defined by
	\[\partial_j(f)=\frac{f-(s_j\cdot f)}{x_j-x_{j+1}}
	=\frac{f(x_1,\ldots,x_n)-f(x_1,\ldots,x_{j-1},x_{j+1},x_j,x_{j+2},\ldots,x_n)}{x_j-x_{j+1}}.\]
	The \emph{K-divided difference operators} $\overline{\partial}_j$ are defined on the polynomial ring $\mathbb{Z}[x_1,\ldots,x_n]$ by
	$\overline{\partial}_j(f)=\partial_j(f-x_{j+1}f)$.
\end{definition}

\begin{definition}
	The \emph{Schubert polynomial} $\mathfrak{S}_w$ of $w\in S_n$ is defined recursively on the weak Bruhat order. Let $w_0=n \hspace{.1cm} n\!-\!1 \hspace{.1cm} \cdots \hspace{.1cm} 2 \hspace{.1cm} 1 \in S_n$, the longest permutation in $S_n$. If $w\neq w_0$ then there is $j\in [n-1]$ with $w(j)<w(j+1)$ (called an \emph{ascent} of~$w$). The polynomial $\mathfrak{S}_w$ is defined by
	\begin{align*}
		\mathfrak{S}_w=\begin{cases}
			x_1^{n-1}x_2^{n-2}\cdots x_{n-1}&\mbox{ if } w=w_0,\\
			\partial_j \mathfrak{S}_{ws_j} &\mbox{ if } w(j)<w(j+1).
		\end{cases}
	\end{align*}
\end{definition}

\begin{definition}
	The \emph{Grothendieck polynomial} $\mathfrak{G}_w$ of $w\in S_n$ is defined analogously to the Schubert polynomial,
	with
	\begin{align*}
		\mathfrak{G}_w=\begin{cases}
			x_1^{n-1}x_2^{n-2}\cdots x_{n-1}&\mbox{ if } w=w_0,\\
			\overline{\partial}_j \mathfrak{G}_{ws_j} &\mbox{ if } w(j)<w(j+1).
		\end{cases}
	\end{align*}	
\end{definition}

Recall that $\ell(w)=\#\{(i,j)\mid 1\leq i<j\leq n \mbox{ with } w(i)>w(j) \}$ is the number of inversions of $w\in S_n$. It can be seen from the recursive definitions above that $\mathfrak{S}_w$ is homogeneous of degree $\ell(w)$, and equals the lowest-degree nonzero homogeneous component of $\mathfrak{G}_w$. See~\cite{manivel} for a deeper introduction to Schubert polynomials. 

\setcounter{section}{2}\setcounter{theorem}{3}
\begin{definition}
	The \emph{Rothe diagram} of a permutation $w\in S_n$ is the set
	\[D(w)=\{(i,j):\, 1\leq i,j\leq n,\,\, w(i)>j,\,\, \mbox{and } w^{-1}(j)>i \}.\]
\end{definition}

We will make use of the following graphical interpretation of $D(w)$. Start with an $n\times n$ grid of boxes with the usual matrix indexing. Place a dot in box $(i,w(i))$ for each $i\in [n]$. Draw rays emanating south and east of each dot, called \emph{hooks}. Remove any boxes hit by a hook. The remaining boxes lie exactly in the positions $D(w)$.

\begin{definition}
	The \emph{Lehmer code} of $w\in S_n$ is the vector
	\[\cc(w)=(\cc(w)_1,\ldots,\cc(w)_n),\]
	where $\cc(w)_i$ equals the number of boxes in row~$i$ of $D(w)$. Symbolically,
	\[\cc(w)_i = \#\{j\mid i<j\mbox{ and } w(j)<w(i) \}. \]
	
\end{definition}

\setcounter{section}{3}\setcounter{theorem}{0}
\begin{definition}\label{def:1}
	An \emph{increasing subsequence} of a permutation $w\in S_n$ starting from position $q\in [n]$ is a vector $\alpha=(\alpha_1,\ldots,\alpha_m)$ such that
	\[q=\alpha_1<\alpha_2<\cdots<\alpha_m\leq n \quad \mbox{and}\quad w(\alpha_1)<w(\alpha_2)<\cdots<w(\alpha_m).\] 
	An increasing subsequence of a permutation~$w$ starting from position~$q$ is called \emph{longest} if there is no longer such sequence, i.e.~$m$ is maximal. We refer to~$m$ as the \emph{length} of~$\alpha$. We denote by $\LIS(w,q)$ the set of all longest increasing subsequences of~$w$ starting from position~$q$.
\end{definition}

\setcounter{section}{3}\setcounter{theorem}{1}
\begin{definition}
	To each $w\in S_n$, associate the vector 
	\[\orr(w)=(\orr(w)_1,\ldots,\orr(w)_n)\in\mathbb{Z}^n,\]
	where for each~$i$, $\orr(w)_i$ equals the length of any $\alpha\in \LIS(w,i)$ minus one. 
\end{definition}

\setcounter{section}{3}\setcounter{theorem}{2}
\begin{definition}[{\cite{PSW}}]
	The \emph{Rajchgot code} of $w\in S_n$ is the vector $\rr(w)\in \mathbb{Z}^n$ with 
	\[\rr(w)_k = n-k - \orr(w)_k \quad\mbox{for each } k\in [n]. \]
\end{definition}

\setcounter{section}{3}\setcounter{theorem}{3}
\begin{definition}
	The \emph{Rajchgot index} of $w\in S_n$ is
	\[\absrajcode{w} = |\rr(w)| = \sum_{k=1}^n \rr(w)_k.\]
	Similarly, let $\absoverrajcode{w} = |\orr(w)|$.
\end{definition}

In general, note that $\rr(w)_k\geq \cc(w)_k$ for all $k\in [n]$. This follows since $\cc(w)_k$ counts the number of inversions of the form $(k,\star)$ in~$w$, and $\rr(w)_k$ counts the number of entries in $[k+1,n]$ excluded in any (fixed) sequence $\alpha\in \LIS(w,k)$.

\section{Climbing chains and markings}
We recall climbing chains, a combinatorial model for Schubert and Grothendieck polynomials due to Lenart, Robinson, and Sottile in~\cite{LRS}. We depict their construction in terms of the Rothe diagrams; this diverges from their exposition, but we find this rendering of climbing chains particularly helpful. 

The \emph{(strong) Bruhat order} is the partial order $\leq$ on $S_n$ defined as the transitive closure of the relations 
\[w\leq wt_{ij} \]
for any $i<j$ such that $w(i)<w(j)$.
We write $w\lessdot v$ to denote a cover relation in the Bruhat order. The following well-known lemma describes all cover relations of the Bruhat order. See for instance~\cite[Lemma 2.1.4]{combcoxgroups} for a proof.

\setcounter{section}{4}\setcounter{theorem}{0}
\begin{lemma}
	\label{lem:cover}
	For $v,w\in S_n$, $w\lessdot v$ if and only if there is $i,j\in [n]$ with $i<j$ such that:
	\begin{itemize}
		\item $v=wt_{ij}$,
		\item $w(i)<w(j)$,
		\item $\{k \mid i<k<j\}\cap \{k\mid w(i)<w(k)<w(j) \}=\emptyset$.
	\end{itemize}
\end{lemma}

Lemma~\ref{lem:cover} can be interpreted graphically as follows. For $i<j$, $w\lessdot wt_{ij}$ if and only if inside $D(w)$:
\begin{itemize}
	\item The dot at $(i,w(i))$ lies north and west of the dot at $(j,w(j))$,
	\item No other dot lies in the region bounded by the hooks at $(i,w(i))$ and $(j,w(j))$.
\end{itemize}

\setcounter{section}{4}\setcounter{theorem}{2}
\begin{definition}
	A \emph{climbing chain} of $w\in S_n$ is a sequence~$C$ of pairs 
	\[(i_1,j_1),\ldots,(i_m,j_m)\]
	called \emph{links}, such that 
	\begin{itemize}
		\item $i_p<j_p$ for each $p\in [m]$, 
		\item $i_1\leq i_2\leq \cdots \leq i_m$,
		\item $wt_{i_1j_1}\cdots t_{i_mj_m} = w_0$,
		\item $wt_{i_1j_1}\cdots t_{i_pj_p} \lessdot wt_{i_1j_1}\cdots t_{i_{p+1}j_{p+1}}$ for each $p\in [m]$ (in the Bruhat order on~$S_n$).
	\end{itemize}
	We call $\ell(C)=m$ the \emph{length} of~$C$, and will write $C_p$ to indicate the link $(i_p,j_p)$ for each $p\in [m]$.
\end{definition}

Observe that all climbing chains of a given permutation~$w$ have the same length $\ell(w_0)-\ell(w)$.

\setcounter{section}{4}\setcounter{theorem}{4}
\begin{definition}
	Associated to each climbing chain~$C$ of a permutation $w\in S_n$ is a set of special links $M(C)$ called the \emph{minimal markings} of~$C$. As~$C$ is a sequence of distinct pairs $(i_p,j_p)$, we will abuse notation slightly and write $M(C)\subseteq C$. If the links of~$C$ are $C_p=(i_p,j_p)$ for each $p\in [\ell(C)]$, then (taking $i_0=0$)
	\[M(C)=\{C_p \mid i_{p-1}<i_{p} \mbox{ or } i_{p-1}=i_p \mbox{ and } j_{p-1}<j_p\}. \]
\end{definition}

Observe that in the trivial case $w=w_0$, the only climbing chain is the empty sequence. We take the empty sequence to have minimal markings $\emptyset$.

\setcounter{section}{4}\setcounter{theorem}{6}
\begin{definition}
	For $w\in S_n$, a \emph{marked climbing chain} of~$w$ is a pair $(C,U)$ where~$C$ is a climbing chain of~$w$ and~$U$ is a superset of the minimal markings of~$C$, that is $M(C)\subseteq U\subseteq C$.
\end{definition}

\setcounter{section}{4}\setcounter{theorem}{7}
\begin{definition}
	The \emph{dual weight} of a marked climbing chain $\xi=(C,U)$ of $w\in S_n$ is the vector \[\dwt(\xi)=(\dwt(\xi)_1,\ldots,\dwt(\xi)_n) \]
	where $\dwt(\xi)_k$ equals the number of links $(k,\star)\in U$.
\end{definition}

\setcounter{section}{4}\setcounter{theorem}{8}
\begin{definition}
	The \emph{weight} of a marked climbing chain $\xi=(C,U)$ of $w\in S_n$ is the vector \[\wt(\xi)=(\wt(\xi)_1,\ldots,\wt(\xi)_n)\] whose~$k$th component equals $n-k-\dwt(\xi)_k$.
\end{definition}

\setcounter{section}{4}\setcounter{theorem}{9}
\begin{theorem}[{\cite[Theorem 5.6]{LRS}}]
	\label{thm:grothformula}
	For any $w\in S_n$,
	\[\mathfrak{G}_w(x_1,\ldots,x_n) = \sum_{\xi} \mathrm{sign}(\xi) \bm{x}^{\wt(\xi)} =\sum_{\xi} \mathrm{sign}(\xi) \frac{\mathfrak{G}_{w_0}}{\bm{x}^{\dwt(\xi)}}, \]
	where the sum is over all marked climbing chains $\xi=(C,U)$ of~$w$, and $\mathrm{sign}(\xi) = (-1)^{\ell(C)-\#U}$.
\end{theorem}

\section{The greedy chain}
\setcounter{section}{5}\setcounter{theorem}{0}
\begin{definition}
	To each permutation $w\in S_n$ with $w\neq w_0$, we associate a pair of integers $\greedy(w)=(a,b)$, where
	\[a=\min\{k\mid w(k)\neq n+1-k\}\quad \mbox{and}\quad b=w^{-1}(w(a)+1).\] 
\end{definition}

\setcounter{section}{5}\setcounter{theorem}{1}
\begin{definition}
	\label{def:greedy}
	Define the \emph{greedy chain} $\CG(w)$ of $w\in S_n$ as follows. If $w=w_0$, then $\CG(w)$ is the empty sequence. Otherwise, define $\CG(w)$ inductively by prepending $\greedy(w)$ to $\CG(wt_{ab})$.
\end{definition}

\section{Algorithm 1 and its properties (original Section 8 core)}
Let~$C$ be a climbing chain of $w\in S_n$. Suppose~$C$ consists of
\[C: w=w^{(0)} \xrightarrow{C_1=(i_1,j_1)} w^{(1)} \xrightarrow{C_2=(i_2,j_2)} w^{(2)}\xrightarrow{C_3=(i_3,j_3)}\cdots\xrightarrow{C_m=(i_m,j_m)} w^{(m)} = w_0. \]
Using Algorithm~\ref{alg:1} below, we associate to~$C$ two sequences of sets: $\Psi_0,\Psi_1,\ldots,\Psi_m$ and $\Omega_0,\Omega_1,\ldots,\Omega_m$.

\begin{algorithm}
	\caption{}
	\label{alg:1}
	\begin{algorithmic}
		\State input $w\in S_n$
		\State input a climbing chain~$C$ of~$w$ with components
		\[C: w=w^{(0)} \xrightarrow{C_1=(i_1,j_1)} w^{(1)} \xrightarrow{C_2=(i_2,j_2)} w^{(2)}\xrightarrow{C_3=(i_3,j_3)}\cdots\xrightarrow{C_m=(i_m,j_m)} w^{(m)} = w_0\]
		
		\\
		\For{$q\in[n-1]$}
		\State compute the lexicographically last element $\alpha^{q}=(\alpha^q_1,\ldots,\alpha^q_{k_q})$ of $\LIS(w,q)$ 
		\EndFor
		\State initialize $\Psi_0 = \bigcup_{q=1}^{n-1}\left\{(q,\alpha^q_2),\ldots,(q,\alpha^q_{k_q}) \right\}$
		\State initialize $\Omega_0=\emptyset$
		\\
		
		\For{$k=1,2,\ldots,m$}
		\If{$C_k\in \Psi_{k-1}$} 
		\State set $\Psi_k=\Psi_{k-1}\setminus\{C_k\}$
		\State set $\Omega_k = \Omega_{k-1}\cup\{C_k\}$
		\ElsIf{$w^{(k)}(a')>w^{(k)}(b')$ for some $(a',b')\in \Psi_{k-1}$}
		\State set $\Psi_k=\{(t_{i_kj_k}(a), b) \mid (a,b)\in \Psi_{k-1} \}$
		\State set $\Omega_k=\Omega_{k-1}$
		\Else
		\State set $\Psi_k=\Psi_{k-1}$
		\State set $\Omega_k = \Omega_{k-1}$
		\EndIf
		\EndFor\\
		\Return $\Psi_0,\ldots,\Psi_m$ and $\Omega_0,\ldots,\Omega_m$
	\end{algorithmic}
\end{algorithm}


\setcounter{section}{8}\setcounter{theorem}{0}
\begin{definition}
	Observe that Algorithm~\ref{alg:1} builds $\Omega_k$ and $\Psi_k$ by performing exactly one of three available transformations on $\Omega_{k-1}$ and $\Psi_{k-1}$: the ``if'', ``else if'', and ``else'' blocks. We will (respectively) name these three operations \emph{transfer}, \emph{adjust}, and \emph{pass}.
\end{definition}

\setcounter{section}{8}\setcounter{theorem}{2}
\begin{lemma}
	\label{lem:algkeyproperty}
	Let $w\in S_n$ and~$C$ be a climbing chain of~$w$. Assume the notation of Algorithm~\ref{alg:1}. Fix~$k$ with $0\leq k\leq \ell(C)$. For each $q\in [n-1]$, let the elements of $\Psi_k$ of the form $(q,\star)$ be $(q,\beta^q_1),\ldots,(q,\beta^q_{l_q})$, labeled so $\beta^q_1<\cdots<\beta^q_{l_q}$. Then
	\[q<\beta^q_1, \quad\mbox{and}\quad w^{(k)}(q) < w^{(k)}(\beta^q_1) < \cdots < w^{(k)}(\beta^q_{l_q}).\]
\end{lemma}

\begin{proof}
	We work by induction on~$k$. When $k=0$, the lemma follows from the definition of $\Psi_0$. Suppose the lemma holds for $k-1$. 
	
	Suppose first that $\Psi_k$ is obtained from $\Psi_{k-1}$ by a transfer operation. Then $\Psi_k=\Psi_{k-1}\setminus \{(i_k,j_k)\}$. Since $(i_k,j_k)$ is a link in~$C$, this means it must be the lexicographically first element of $\Psi_{k-1}$ (by the induction assumption). The statement of the lemma follows easily in this case.
	
	Next, suppose that $\Psi_k$ is obtained from $\Psi_{k-1}$ by an adjust operation (see Figure~\href{https://alco.centre-mersenne.org/articles/10.5802/alco.358/}{4} for a visual representation of this case). It follows $(i_k,j_k)\notin \Psi_{k-1}$, $w^{(k)}(a')>w^{(k)}(b')$ for some $(a',b')\in \Psi_{k-1}$, and \[\Psi_k = \{(t_{i_kj_k}(a),b)\mid (a,b)\in \Psi_{k-1} \}. \] Observe that all such $(a',b')$ must satisfy $a'=i_k$, $b'>j_k$, and $w^{(k-1)}(i_k)<w^{(k-1)}(b')<w^{(k-1)}(j_k)$. The adjust operation then guarantees the conditions of the lemma are met.
	
	Lastly, consider the case of a pass operation. This means that $(i_k,j_k)\notin \Psi_{k-1}$ and $w^{(k)}(a)<w^{(k)}(b)$ for all $(a,b)\in \Psi_{k-1}$. The conditions of the lemma are immediate.
\end{proof}

\setcounter{section}{8}\setcounter{theorem}{3}
\begin{lemma}
	\label{lem:adjustcase}
	Let $w\in S_n$ and~$C$ be a climbing chain of~$w$. Assume the notation of Algorithm~\ref{alg:1}. Suppose step~$k$ in the execution of Algorithm~\ref{alg:1} is an adjust operation caused by $(a',b')\in \Psi_{k-1}$. Then $a'=i_k$ and $b'>j_k$.
\end{lemma}

\setcounter{section}{8}\setcounter{theorem}{4}
\begin{lemma}
	\label{lem:algbasicproperties}
	Let $w\in S_n$ and~$C$ be a climbing chain of~$w$. Assume the notation of Algorithm~\ref{alg:1}. Then 
	\begin{enumerate}[label=\textup{(\roman*)}]
		\item $\emptyset=\Omega_0\subseteq \cdots\subseteq \Omega_m\subseteq \{C_1,\ldots,C_m\}$.
		\item For each $0\leq k\leq m$, one has $\Omega_k\subseteq \{C_1,\ldots,C_k \}$.
		\item If $(i_p,j_p)\in \Omega_k$ for some~$k$, then $\Omega_p=\Omega_{p-1}\cup \{(i_p,j_p)\}$ and $\Psi_p=\Psi_{p-1}\setminus \{(i_p,j_p)\}$.
		\item For each $0\leq k\leq m$, one has $\#\Omega_k + \#\Psi_k = \absoverrajcode{w}.$
		\item $\Psi_m=\emptyset$ and $\#\Omega_m=\absoverrajcode{w}$.
	\end{enumerate}
\end{lemma}

\begin{proof}
	We first address (i) and (ii). Initially, $\Omega_0=\emptyset$. Only the transfer operation causes $\Omega_k\neq \Omega_{k-1}$, specifically by adding the element $C_k=(i_k,j_k)$ to $\Omega_{k-1}$. Thus~(i) and (ii) hold. For (iii), note that $(i_p,j_p)\in \Omega_k$ implies step~$p$ of Algorithm~\ref{alg:1} was a transfer operation. Claim (iv) holds by an argument analogous to that of (ii). The claim that $\Psi_m=\emptyset$ in (v) follows from Lemma~\ref{lem:algkeyproperty}. That $\#\Omega_m=\absoverrajcode{w}$ then follows from (iv).
\end{proof}

\setcounter{section}{8}\setcounter{theorem}{5}
\begin{definition}
	Let~$C$ be a climbing chain of $w\in S_n$. Suppose the minimal markings $M(C)$ are 
	\[\{C_{k_1},C_{k_2},\ldots,C_{k_p}\}\quad \mbox{with}\quad 1\leq k_1<k_2<\cdots<k_p\leq \ell(C).\]
	Define the \emph{runs} of~$C$ to be the (disjoint) subsequences of~$C$ consisting of links $k_q$ through $k_{q+1}-1$ for each $q\in [p]$ (taking $k_{p+1}=\ell(C)+1$).
\end{definition}

\setcounter{section}{8}\setcounter{theorem}{7}
\begin{lemma}
	\label{lem:atmostoneperrun}
	Let $w\in S_n$ and~$C$ be a climbing chain of~$w$ with length~$m$. Construct the sets $\Omega_0,\ldots,\Omega_m$ and $\Psi_0,\ldots,\Psi_m$ using Algorithm~\ref{alg:1}. Then each run of~$C$ contains at most one element of $\Omega_m$.
\end{lemma}

\begin{proof}
	Consider a run $C_{p},C_{p+1},\ldots,C_{q}$ of~$C$ containing at least two elements of $\Omega_m$. The definition of run forces $\{C_p,C_{p+1},\ldots,C_q\}\cap M(C) = \{C_p\}$. It follows that $i_p=i_{p+1}=\cdots=i_{q}$ and $j_p>j_{p+1}>\cdots>j_q>i_p$. 
	
	First, suppose that $C_a,C_{a+1}\in \Omega_m$ for some $p\leq a< q$. By Lemma~\ref{lem:algbasicproperties}(iii), $C_a,C_{a+1}\in \Psi_{a-1}$. But $i_a=i_{a+1}$ and $j_a>j_{a+1}$, so Lemma~\ref{lem:algkeyproperty} implies
	\[w^{(a-1)}(i_a) < w^{(a-1)}(j_{a+1}) < w^{(a-1)}(j_a). \]
	This contradicts the climbing chain condition $w^{(a-1)}\lessdot w^{(a)} = w^{(a-1)}t_{i_aj_a}$ (see Lemma~\ref{lem:cover}).
	
	Now suppose $C_a,C_b\in \Omega_m$ for some $a,b$ with $p\leq a<b-1$, $b\leq q$ and $b-a$ minimal. Additionally, suppose there are no adjust operations between steps~$a$ and~$b$ in the execution of Algorithm~\ref{alg:1}. This implies $\Psi_a=\Psi_{a+1}=\cdots=\Psi_{b-1}$. Thus $C_a,C_b\in \Psi_{a-1}$, again contradicting the Bruhat cover condition on climbing chains.
	
	Lastly, suppose there was an adjust operation between steps~$a$ and~$b$ in the execution of Algorithm~\ref{alg:1}. Say the first such adjust operation occurs at step~$k$. Recall that $i_a=i_{a+1}=\cdots=i_k=\cdots=i_b$. Suppose first that there are no elements $(i_k,q)\in \Psi_{k}$. Then by Lemma~\ref{lem:adjustcase}, no further adjust operations occur prior to step~$b$. Thus there are no elements $(i_k,q)\in \Psi_{b-1}$. By Lemma~\ref{lem:algbasicproperties}, this contradicts that $C_b\in \Omega_m$.
	
	Hence we may assume that there is an element of the form $(i_k,q)\in \Psi_k$. Then $(j_k,q)\in \Psi_{k-1}$, so Lemma~\ref{lem:algkeyproperty} implies $q>j_k$. Hence $q>j_k>j_{k+1}>\cdots>j_b$, so $(i_k,q)\notin \{C_k,C_{k+1},\ldots,C_{b}\}$. 
	
	If there is only a single adjust operation between steps~$a$ and~$b$, this contradicts that $C_b\in \Omega_m$. If there is more than one adjust operation between steps~$a$ and~$b$, the same contradiction is reached by repeating the previous argument for each adjust operation. 
\end{proof}

\setcounter{section}{8}\setcounter{theorem}{8}
\begin{theorem}
	\label{prop:minnummarks}
	For any $w\in S_n$ and any climbing chain~$C$ of~$w$, 
	\[\absoverrajcode{w} \leq \#M(C). \]
\end{theorem}

\begin{proof}
	Use Algorithm~\ref{alg:1} on~$C$ to produce $\Omega_m$. Decompose~$C$ into runs. By definition, the number of runs equals $\#M(C)$. Lemma~\ref{lem:atmostoneperrun} shows there is at most one element of $\Omega_m$ in each run. Lemma~\ref{lem:algbasicproperties}(v) implies $\#\Omega_m=\absoverrajcode{w}$. Putting this all together,
	\[\#M(C)\geq \#\Omega_m=\absoverrajcode{w}.\qedhere\]
\end{proof}

\section{Leading monomials via heavy climbing chains (original Section 9 core)}
\setcounter{section}{9}\setcounter{theorem}{0}
\begin{lemma}
	\label{lem:drecursionanychain}
	Let~$C$ be a climbing chain of $w\in S_n$ with $w\neq w_0$. Suppose $C_p=(i_p,j_p)$ for $p\in [\ell(C)]$, and set $v=wt_{i_1j_1}$. Then 
	\begin{enumerate}[label=\textup{(\roman*)}]
		\item $\orr(w)_p=0=\orr(v)_p$ for $p\in [i_1-1]$,
		\item $\orr(w)_{i_1}\geq \orr(v)_{i_1}$,
		\item $\orr(w)_p\leq \orr(v)_p$ for $i_1+1\leq p \leq n$.
	\end{enumerate}
\end{lemma}

\begin{proof}
	Claim (i) is immediate since $i_1$ since $w(k)=v(k)=n-k+1$ for $k\in[i_1]$. Since $v(i_1)>w(i_1)$, any $\alpha\in \LIS(v,i_1)$ is also an increasing sequence in~$w$. This shows (ii). 
	
	Consider claim (iii). When $p>j_1$, (iii) follows since~$w$ and~$v$ agree after position $j_1$. Since $v(j_1)<w(j_1)$, (iii) holds when $p=j_1$. It remains to consider $i_1<p<j_1$. Since $w\lessdot v$, either $w(p)<w(i_1)$ or $w(p)>w(j_1)$. In either case, any $\alpha\in \LIS(w,p)$ is also an increasing sequence in~$v$. This proves (iii).
\end{proof}

\setcounter{section}{9}\setcounter{theorem}{1}
\begin{definition}
	We call a climbing chain~$C$ of $w\in S_n$ \emph{heavy} if $\#M(C) = \absoverrajcode{w}$.
\end{definition}
We choose the name ``heavy'' since such chains~$C$ will have maximal total weight: $|\wt(C)| = \absrajcode{w}$.

\setcounter{section}{9}\setcounter{theorem}{2}
\begin{lemma}
	\label{lem:heavytruncation}
	Let~$C$ be a climbing chain of $w\in S_n$, and assume the notation of Algorithm~\ref{alg:1}.
	Suppose $C_k$ is the last link in~$C$ of the form $(i_1,\star)$. Then 
	\begin{itemize}
		\item $\orr(w)_p=0=\orr(w^{(k)})_p$ for $p\in [i_1-1]$,
		\item $\orr(w)_{i_1}\geq \orr(w^{(k)})_{i_1}=0$,
		\item $\orr(w)_p\leq \orr(w^{(k)})_p$ for $i_1+1\leq p \leq n$.
	\end{itemize}
	whenever $i_1+1\leq p\leq n$. Additionally if~$C$ is heavy, then the truncation $C'=(C_{k+1},\ldots,C_m)$ is heavy (as a climbing chain of $w^{(k)}$).
\end{lemma}

\begin{proof}
	The itemized claims follow from repeated applications of Lemma~\ref{lem:drecursionanychain}. To see the final claim, suppose $C'$ is not heavy. Concatenating $(C_1,\ldots,C_{k})$ with a heavy chain for $(w^{(k)})$ yields a climbing chain of~$w$ with fewer minimal marking than~$C$, implying that~$C$ is not heavy.
\end{proof}

\setcounter{section}{9}\setcounter{theorem}{3}
\begin{lemma}
	\label{lem:heavymaxfirstmarks}
	Let~$C$ be a heavy climbing chain of $w\in S_n$. Then
	\[\dwt(C,M(C))_{i_1}\leq \orr(w)_{i_1}. \]
\end{lemma}

\begin{proof}
	Assume the notation of Algorithm~\ref{alg:1}. Suppose~$k$ is the index of the last link in~$C$ of the form $(i_1,\star)$ and $C'=(C_{k+1},\ldots,C_m)$.
	By the choice of~$k$, we have
	\begin{itemize}
		\item $\dwt(C,M(C))_p=\dwt(C',M(C'))_p=0$ for $p<i_1$,
		\item $\dwt(C,M(C))_{i_1} >0$,
		\item $\dwt(C',M(C'))_{i_1} =0$,
		\item $\dwt(C,M(C))_p=\dwt(C',M(C'))_p$ for $p>i_1$.
	\end{itemize}
	Observe that the choice of~$k$ implies $\orr(w^{(k)})_{p}=0$ for $p\leq i_1$. Applying Lemma~\ref{lem:heavytruncation},
	\begin{align*}
		\left|\dwt(C,M(C))\right| &= \sum_{p=1}^{n}\dwt(C,M(C))_p \\
		&= \dwt(C,M(C))_{i_1} + \sum_{p=i_1+1}^{n}\dwt(C,M(C))_{p} \\
		&= \dwt(C,M(C))_{i_1} + \sum_{p=i_1+1}^{n}\dwt(C',M(C'))_{p} \\
		&= \dwt(C,M(C))_{i_1} + \left|\dwt(C',M(C'))\right| \\
		&= \dwt(C,M(C))_{i_1} +\absoverrajcode{w^{(k)}} \\
		&= \dwt(C,M(C))_{i_1} +\orr(w^{(k)})_{i_1+1}+\cdots +\orr(w^{(k)})_n.
	\end{align*}
	On the other hand,~$C$ being heavy implies
	\begin{align*}
		\left|\dwt(C,M(C))\right| &= \absoverrajcode{w}\\
		&= \orr(w)_{i_1}+\cdots+\orr(w)_n.
	\end{align*}
	Hence,
	\begin{align}
		\label{eqn:1}
		\biggr(\dwt(C,M(C))_{i_1}-\orr(w)_{i_1}\biggr)+\sum_{p=i_1+1}^{n} \biggr(\orr(w^{(k)})_{p} -\orr(w)_{p}\biggr)=0.
	\end{align}
	By Lemma~\ref{lem:drecursionanychain}, each term in the summation is nonnegative, so the leftmost term must be nonpositive. Thus,
	\[\dwt(C,M(C))_{i_1}\leq \orr(w)_{i_1}.\qedhere \] 
\end{proof}

\setcounter{section}{9}\setcounter{theorem}{4}
\begin{lemma}
	\label{lem:heavyequalitycase}
	Let~$C$ be a heavy climbing chain of $w\in S_n$. Assume the notation of Algorithm~\ref{alg:1}. 
	Suppose that~$k$ is the index of the last link in~$C$ of the form $(i_1,\star)$, and
	\[\dwt(C,M(C))_{i_1}= \orr(w)_{i_1}. \]
	Then
	\[
	\orr(w^{(k)})_p = \orr(w)_p \quad\mbox{for } p\geq i_1+1. 
	\]
\end{lemma}

\begin{proof}
	This is an immediate consequence of~\eqref{eqn:1}.
\end{proof}

\setcounter{section}{9}\setcounter{theorem}{5}
\begin{theorem}
	\label{thm:leadingchainweight}
	For $w\in S_n$ and any term order satisfying $x_1<x_2<\cdots<x_n$,
	\[\bm{x}^{\orr(w)} = \min \left\{\bm{x}^{\dwt(C,M(C))} \mid C \mbox{ is a heavy climbing chain of } w \right\}. \] 
	Equivalently,
	\[\bm{x}^{\rr(w)} = \max \left\{\bm{x}^{\wt(C,M(C))} \mid C \mbox{ is a heavy climbing chain of } w \right\}. \] 
\end{theorem}

\begin{proof}
	It is immediate from the definition of $\wt$ that the two assertions of the theorem are equivalent. We focus on the first. Since all heavy climbing chains have the same number of minimal markings, the theorem is equivalent to proving 
	\[\bm{x}^{\orr(w)} = \max \left\{\bm{x}^{\dwt(C,M(C))} \mid C \mbox{ is a heavy climbing chain of } w \right\}. \] 
	in any term order with $x_1>x_2>\cdots>x_n$. 
	For each $p\in[n]$, let $k_p$ be the index of the last link in~$C$ of the form $(p,\star)$. The theorem follows by applying Lemmas~\ref{lem:heavymaxfirstmarks} and~\ref{lem:heavyequalitycase} (sequentially) to each of $w^{(k_1)},w^{(k_2)},\ldots,w^{(k_{n-1})}$.
\end{proof}

\section{Interpolating chains and the vector sequence $L_w$}
\setcounter{section}{10}\setcounter{theorem}{3}
\begin{definition}
	To each permutation $w\in S_n$, we associate a sequence of vectors \[\leads(w)=(L_w(0),\ldots,L_w(d)),\] where $d=\absrajcode{w}-\absinvcode{w}$. Set $L_w(0)=\cc(w)$. Define $L_w(i)$ for $i\in [d]$ as follows: set $j\in [n]$ to be the largest index such that $L_w(i-1)_j<\rr(w)_j$, and define $L_w(i) = L_w(i-1) + e_j$.
\end{definition}

Recall that $\rr(w)_k\geq \cc(w)_k$ for all $k\in [n]$, so $L_w(d) = \rr(w)$. 

\section{Proof of the Hafner conjecture: leading monomials above minimal degree}
In Corollary~\href{https://alco.centre-mersenne.org/articles/10.5802/alco.358/}{10.10} we showed that the monomials named in Conjecture~\ref{conj:elena} are in the support of $\mathfrak{G}_w$. We show that these monomials are actually the leading monomials of the homogeneous components of $\mathfrak{G}_w$.

\setcounter{section}{11}\setcounter{theorem}{0}
\begin{lemma}
	\label{lem:maxmarkedones}
	Let $(C,U)$ be any marked climbing chain of $w\in S_n$ and $C_1=(i_1,j_1)$. Then 
	\[\dwt(C,U)_{i_1} \leq \occ(w)_{i_1}.\]
	If equality holds, then 
	\[(C_1,C_2,\ldots,C_k) = (\CG_1,\ldots,\CG_k), \]
	where $C_k$ is the last link in~$C$ of the form $(i_1,\star)$.
\end{lemma}

\begin{proof}
	From the definition of a climbing chain, we see that all links of the form $(i_1,\star)$ in~$C$ are noninversions of~$w$, pairs $(a,b)$ with $a<b$ and $w(a)<w(b)$. There are exactly $n-i_1-\cc(w)_{i_1} = \occ(w)_{i_1}$ of these noninversions with $a=i_1$. The second assertion follows from the nested hook interpretation of climbing chain links.
\end{proof}

\setcounter{section}{11}\setcounter{theorem}{1}
\begin{lemma}
	\label{lem:onemarkbasecase}
	Let $(C,U)$ be a marked climbing chain of $w\in S_n$ with~$C$ given by
	\[C: w=w^{(0)} \xrightarrow{C_1=(i_1,j_1)} w^{(1)} \xrightarrow{C_2=(i_2,j_2)} w^{(2)}\xrightarrow{C_3=(i_3,j_3)}\cdots\xrightarrow{C_m=(i_m,j_m)} w^{(m)} = w_0.\]
	Then
	\[\dwt(C,U)_{i_1} \leq \LL_w(m-\#U)_{i_1}. \]
\end{lemma}

\begin{proof}
	Suppose $C_k$ is the last link in~$C$ of the form $(i_1,\star)$. 
	By Theorem~\ref{prop:minnummarks},
	\[\#U = |\dwt(C,U)| = \dwt(C,U)_{i_1}+\sum_{p=i_1+1}^n\dwt(C,U)_p \geq \dwt(C,U)_{i_1} + \absoverrajcode{w^{(k)}}. \]
	On the other hand, Lemma~\ref{lem:heavytruncation} implies 
	\[\absoverrajcode{w} \leq \orr(w)_{i_1} + \absoverrajcode{w^{(k)}},\]
	so that
	\[\#U = \absoverrajcode{w} + (\#U - \absoverrajcode{w})\leq \orr(w)_{i_1} + \absoverrajcode{w^{(k)}} + (\#U - \absoverrajcode{w}).\]
	Hence
	\[\dwt(C,U)_{i_1} \leq \orr(w)_{i_1} + (\#U - \absoverrajcode{w}).\]
	
	To conclude the proof, we will analyze $\orr(w)_{i_1} + (\#U - \absoverrajcode{w})$.
	Let $d=m-\#U$. See Figure~\href{https://alco.centre-mersenne.org/articles/10.5802/alco.358/}{7} for a visual representation of the following argument. Recall that $L_w(d)$ is obtained from $\cc(w)$ by iteratively increasing components~$d$ times, moving right-to-left so entries stay weakly below the corresponding entries of $\rr(w)$. Then $\LL_w(d)$ is obtained from $\occ(w)$ by iteratively decreasing components~$d$ times, moving right-to-left so entries stay weakly above the corresponding entries of $\orr(w)$.
	
	Equivalently, $\LL_w(d)$ is obtained from $\orr(w)$ by iteratively increasing components $m-\absoverrajcode{w}-d$ times, moving left-to-right so entries stay weakly below the corresponding entries of $\occ(w)$. Note that 
	\[m-\absoverrajcode{w}-d = \#U-\absoverrajcode{w}.\]
	If $\orr(w)_{i_1} + (\#U - \absoverrajcode{w})\geq \occ(w)_{i_1}$, then $\LL_w(d)_{i_1}=\occ(w)_{i_1}$ and Lemma~\ref{lem:maxmarkedones} completes the proof. 
	
	Otherwise, $\orr(w)_{i_1} + (\#U - \absoverrajcode{w})< \occ(w)_{i_1}$, so 
	\[\LL_w(m-\#U)_{i_1} = \orr(w)_{i_1} + (\#U - \absoverrajcode{w}) \geq \dwt(C,U)_{i_1}\]
	as shown above, so we are done.
\end{proof}

\setcounter{section}{11}\setcounter{theorem}{2}
\begin{theorem}
	\label{thm:interpolatingleadingterms}
	Fix $w\in S_n$, and let $\mathcal{Z}_w$ denote the set of marked climbing chains of $w$. For each~$d$ with $1\leq d\leq \absrajcode{w}-\absinvcode{w}$ and any term order satisfying $x_1<x_2<\cdots<x_n$,
	\[\bm{x}^{\LL_w(d)} = \min \left\{\bm{x}^{\dwt(C,U)} \mid (C,U)\in \mathcal{Z}_w \mbox{ with } \#U = 
	\ell(C) - d \right\}. \] 
	Equivalently, 
	\[\bm{x}^{L_w(d)} = \max \left\{\bm{x}^{\wt(C,U)} \mid (C,U)\in \mathcal{Z}_w \mbox{ with } \#U = 
	\ell(C) - d \right\}. \] 
\end{theorem}

\begin{proof}
	It is immediate from the definition of $\wt$ that the two assertions of the theorem are equivalent. We focus on the first. Since we are considering marked climbing chains with a fixed number of marks, the theorem is equivalent to proving
	\[\bm{x}^{\LL_w(d)} = \max \left\{\bm{x}^{\dwt(C,M(C))} \mid (C,U)\in \mathcal{Z}_w \mbox{ with } \#U = 
	\ell(C) - d \right\}. \] 
	in any term order with $x_1>x_2>\cdots>x_n$. 
	
	Assume the notation of Algorithm~\ref{alg:1}. 
	By definition, $\LL_w(d)_{p}=\dwt(C,U)_{p}$ for $p<i_1$. Lemma~\ref{lem:onemarkbasecase} shows that $\LL_w(d)_{i_1}\geq\dwt(C,U)_{i_1}$. If the inequality is strict, then there is nothing to prove. 
	
	Suppose first that $\LL_w(d)_{i_1}=\dwt(C,U)_{i_1}<\occ(w)_{i_1}$. Then for $p>i_1$, $\LL_w(d)_p=\orr(w)_p$. 
	Let $C_k$ be the last link in~$C$ of the form $(i_1,\star)$. Let $C'=(C_{k+1},\ldots,C_m)$ and $U'=U\cap C'$.
	
	By Theorem~\ref{prop:minnummarks},
	\[\sum_{p=i_1+1}^n \dwt(C,U)_p = |\dwt(C',U')|\geq \absoverrajcode{w^{(k)}}. \]
	For $p>i_1$, Lemma~\ref{lem:heavytruncation} shows $\orr(w^{(k)})_p\geq \orr(w)_p$. Then
	\begin{align}
		\label{eqn:2}
		\sum_{p=i_1+1}^n \dwt(C,U)_p \geq \absoverrajcode{w^{(k)}}\geq \sum_{p=i_1+1}^n\orr(w)_p = \sum_{p=i_1+1}^n\LL_w(d)_p.
	\end{align}
	
	Since $|\LL_w(d)| = |\dwt(C,U)|$ and $\LL_w(d)_p = \dwt(C,U)_p$ for $p\leq i_1$, it follows that equality holds throughout in~\eqref{eqn:2}. Thus $C'$ is heavy for $w^{(k)}$, so the theorem now follows from Theorem~\ref{thm:leadingchainweight}.
	
	It remains to consider the case that $\LL_w(d)_{i_1}=\dwt(C,U)_{i_1}=\occ(w)_{i_1}$. In this case, we may pass to $w^{(k)}$ and repeat the above arguments and iterate. The iteration will terminate at worst with $w_0$, for which the theorem is trivial.
\end{proof}

\begin{thebibliography}{99}
\bibitem[21]{PSW} Oliver Pechenik, David E. Speyer, and Anna Weigandt, ``Regularity of matrix Schubert varieties'', \emph{S\'em. Lothar. Combin.}\textbf{86B}(2022), article no.47 (12pages).
\bibitem[8]{elena1} Elena S. Hafner, ``Vexillary Grothendieck polynomials via bumpless pipe dreams'',2022, arXiv:2201.12432.
\bibitem[17]{manivel} Laurent Manivel, \emph{Symmetric functions, Schubert polynomials and degeneracy loci}, SMF/AMS Texts and Monographs6, American Mathematical Society; Soci\'et\'e Math\'ematique de France,2001, translated from the1998 French original by John R. Swallow.
\bibitem[16]{LRS} Cristian Lenart, Shawn Robinson, and Frank Sottile, ``Grothendieck polynomials via permutation patterns and chains in the Bruhat order'', \emph{Amer. J. Math.}\textbf{128}(2006),no.4,805--848.
\bibitem[3]{combcoxgroups} Anders Bj\"orner and Francesco Brenti, \emph{Combinatorics of Coxeter groups}, Graduate Texts in Mathematics231, Springer, New York,2005.
\end{thebibliography}
\end{document}
